\section{Code and data}
If we want to follow a calculation into the implementation, the table below gives the starting
points in the companion repository. Paths are relative to its root. The full implementations,
saved measurements, and longer derivations are there too: the detailed report remains in
\path{reports/search-math/}, and this edition is in \path{reports/search-math-concise/}.

\begin{center}\small
\begin{tabularx}{\linewidth}{@{}lY@{}}\toprule
Question & Start here\\\midrule
Follow the six-document example & \path{examples/small_search.py}\\
Inspect scan and branching & \path{cpp/index.cpp}, \path{cpp/score.hpp}\\
Inspect key lookup & \path{cpp/key_index.cpp}\\
Check storage and work formulas & \path{experiments/models.py}, \path{experiments/check_controlled.py}\\
Read the final timing checks & \path{experiments/check_native_model.py}, \path{reports/search-math/evidence/final-study.json}\\
Calculate the larger cases & \path{experiments/project_costs.py}, \path{results/projected-cases-2026-09-11/}\\
Reproduce the study & \path{experiments/README.md}, \path{results/README.md}\\\bottomrule
\end{tabularx}\normalsize\end{center}

Once the project is installed using its root README, we can run the six-document example
without downloading a corpus. The second command checks the long report's embedded arithmetic
and saved evidence. The third uses frozen inputs to recalculate the larger scenarios, so it
doesn't build a billion-row index.
\Needspace{12\baselineskip}
\begin{lstlisting}[caption={Three existing entry points, run from the repository root.}]
.venv/bin/python examples/small_search.py
python3 reports/search-math/verify_math.py
.venv/bin/python -m experiments.project_costs \
  --inputs results/projected-cases-2026-09-11/results.json \
  --output results/my-scenarios \
  --documents 1000000000 --ram-gb 32 64 1000
\end{lstlisting}
For the exact saved-input invocation and full-data runs, use the reproduction README. A full
rerun also needs the dataset, embedding cache, and documented environment. Compiling this PDF
only rebuilds the explanation: its tables summarize saved results, not new measurements.

\begin{thebibliography}{9}\small\raggedright
\bibitem{exa} Exa Team. \emph{How we built a web-scale vector database}. 2024.
\url{https://exa.ai/blog/building-web-scale-vector-db}.
\bibitem{mrl} A. Kusupati et al. \emph{Matryoshka Representation Learning}. 2022.
\url{https://arxiv.org/abs/2205.13147}.
\bibitem{msmarco} Microsoft. \emph{MS MARCO}. Passage collection and benchmark resources.
\url{https://microsoft.github.io/msmarco/}.
\bibitem{nomic} Nomic AI. \emph{nomic-embed-text-v1.5}. Model documentation.
\url{https://huggingface.co/nomic-ai/nomic-embed-text-v1.5}.
\bibitem{local} Local experimental record, 10--11 September 2026. Source checkpoint
\texttt{f16ca66}; \path{reports/search-math/evidence/} and the result files it identifies.
The concise edition's README records the selected sources and reproduction scope.
\end{thebibliography}
